\documentclass[11pt]{article}
\usepackage[utf8]{inputenc}
\usepackage[T1]{fontenc}
\usepackage{lmodern}
\usepackage[margin=1in]{geometry}
\usepackage{amsmath,amssymb}
\usepackage{textcomp}
\usepackage[dvipsnames]{xcolor}
\usepackage{enumitem}
\usepackage{parskip}
\usepackage{titlesec}
\usepackage[colorlinks=true,linkcolor=blue,citecolor=blue,urlcolor=blue]{hyperref}
\usepackage[numbers]{natbib}
\usepackage{microtype}
\emergencystretch=1.5em

\titleformat{\section}{\large\bfseries}{\thesection}{1em}{}
\titleformat{\subsection}{\normalsize\bfseries}{\thesubsection}{1em}{}
\titleformat{\subsubsection}{\normalsize\bfseries\itshape}{}{0em}{}

\title{\textbf{Fixed by Construction, Estimated by Generator:\\ The Robustness Limits of Corpus-Level Contamination Estimators}}
\author{}
\date{}

\newcommand{\ph}{\textit{[placeholder --- fill in after experiments are run]}}

\begin{document}
\maketitle
\vspace{-2em}

\begin{quote}
\small\textit{Assembled from the six confirmed stage outputs of an \texttt{auto-research-discussion} run (seed topic B6.1), post research-taste audit (\textcolor{ForestGreen}{\textbf{PASS}}). Status: \textbf{specification / paper draft --- no experiments have been run.} All numbers below are pre-registered design targets, not results.}
\end{quote}

\begin{abstract}
Teams deciding whether to filter or down-weight a training corpus increasingly rely on corpus-level estimators of its synthetic-text fraction, yet these estimators --- including the distributional mixture estimators already deployed to monitor AI-modified content at scale \citep{MonitoringAIContent2024} --- have only ever been validated against a fixed, known set of calibration generators, a condition that fails to hold once a corpus is contaminated by generators the auditor never anticipated. Robustness to unseen generators and to paraphrase attacks has been measured, thoroughly, for \emph{per-document} text detectors \citep{RAID2024Benchmark,MAGE2024InTheWild}; what has not been measured is whether those per-document degradations survive aggregation into a corpus-wide proportion, where an averaging argument predicts they might partially cancel. No prior work reports how a corpus-level estimator's error scales jointly with generator novelty and paraphrase strength, separates that degradation from a held-out generator's intrinsic difficulty, or certifies whether any statistical feature survives both stresses at once. We build a controlled measurement testbed that withholds an entire generator from calibration rather than merely a sample of its output, sweeps paraphrase strength continuously, contrasts a quantification-corrected estimator against a naive one to isolate aggregation bias, and uses a capability-matched seen/unseen swap on identical documents to separate genuine unseen-ness from a generator's intrinsic difficulty. A joint mixed-effects specification disentangles the two stress axes from each other, and a pre-committed cross-cell flatness-and-floor test certifies whether a generator-invariant feature exists. We evaluate this design across controlled human/synthetic mixtures built from a multi-family generator pool and report the resulting degradation law together with a feature-invariance verdict --- a negative result on the latter is a reportable deliverable, not a failure.
\end{abstract}

\section{Introduction}

Corpus-level estimators of synthetic-text contamination have become the practical instrument teams reach for when deciding whether to filter or down-weight a training corpus, and the stakes behind that decision are not cosmetic. Pretraining corpora already contain substantial content produced by many different, often unidentified generating systems, and the fraction of a corpus that is synthetic is not a side statistic: theoretical analyses of iterative retraining on real-plus-synthetic mixtures show that the mixing parameter determines whether repeated retraining on that corpus is a stable process or a collapsing one. Whatever number a contamination estimator returns is, in a very literal sense, an input to whether a downstream training pipeline is stable.

Estimators of this kind are not hypothetical. A distributional maximum-likelihood mixture framework has already been deployed at scale to estimate the fraction of AI-modified content in conference peer-review pools and in the scientific literature, without watermarks and without per-document labels \citep{MonitoringAIContent2024,MappingLLMUse2024} --- the same class of instrument a data-curation team would reach for. Yet a central question about these estimators remains unmeasured: every one of them is calibrated against a fixed, known set of generators and a fixed human reference, and real corpora are contaminated by generators no auditor selected or anticipated in advance.

What is genuinely known about that failure is known one level down, at the per-document detector. Multi-generator benchmarks have measured detector degradation across unseen generator families and across adversarial attacks including paraphrase, in text, at considerable scale \citep{RAID2024Benchmark,MAGE2024InTheWild}; the per-document story is not in dispute and this paper does not re-litigate it. The unmeasured question is what happens one level \emph{up}. Aggregating many per-document readings into a single proportion is not a neutral operation: an averaging argument predicts that independent per-document errors partially cancel in the aggregate, which would make a corpus-level estimator strictly more robust than the detector inside it, while quantification theory predicts the opposite --- that classify-and-count aggregation acquires a systematic bias precisely under distribution shift, in a direction that does not cancel \citep{QuantDatasetShift2023}. These two predictions disagree, and no reported measurement adjudicates between them for a text corpus under an unseen generator.

Prior work approaches the surrounding territory from three directions, none of which reaches this measurement. Zero-shot and rewrite-distance detectors score a single passage against a curvature, distance, or recovery-similarity signal, but their target is a per-document decision, and several require access to a source model's likelihood that a truly unknown generator will not provide; the multi-generator robustness benchmarks that stress them inherit that same per-document estimand. Watermarking sidesteps detection entirely by embedding a signal at generation time, and the sharpest available proportion-estimation theory is built on top of it --- but only for generators that cooperate with the scheme, the complement of a setting where the contaminating generator was never under the auditor's control. Closest in \emph{spirit} is a detector-generalization study in the image modality, which trains a detector to remain accurate on diffusion-model variants unseen during training via contrastive hard-sample augmentation; it is a cross-modal precedent for the failure mode, not the nearest neighbor in text. What none of the three does is measure how a text corpus's estimated \emph{proportion} synthetic degrades as a joint function of an unseen generator and paraphrase strength, or separate that degradation from a held-out generator's intrinsic difficulty --- which is exactly the gap this paper is built to close.

Turning that gap into a measurable question surfaces several distinct difficulties. The space of possible text generators is open-ended and continually growing, so any calibration set is necessarily incomplete, and testing generalization to a \emph{structurally} unseen generator requires withholding an entire generator's documents from calibration rather than merely holding out a sample of a generator already seen. Paraphrasing interpolates continuously between clearly-synthetic and clearly-human text, so the separability margin --- and any single point estimate of contamination --- is inherently unstable near that margin regardless of which estimator is used. Aggregating many per-document readings into one corpus-level ratio is a known source of systematic bias under distribution shift, and that bias moves with the held-out generator rather than washing out with a fixed calibration constant. Even if a leave-one-generator-out gap is observed, it is not obvious whether that gap reflects genuine unseen-ness or simply that the particular held-out generator's text happens to be intrinsically hard to detect --- two explanations that look identical in a plain comparison across different generators. Generator novelty and paraphrase strength also act on the same documents at once, so a test of either factor in isolation risks attributing one axis's effect to the other. And finally, a claim that some statistical feature is ``generator-invariant'' is only meaningful if it is checked in the worst tested condition, not just on average across many --- a feature that separates well in most conditions and collapses in one has not, in the sense that matters for practice, been shown invariant at all.

Our approach turns each of these difficulties into a specific design choice. Because the true contamination ratio is fixed by construction whenever a controlled mixture is built, while every candidate estimator's decision statistic is free to vary with generator identity and paraphrase strength, the central design problem is to isolate exactly how much of that statistic's variation is attributable to each stress axis, separately from ordinary per-document classification error and from a generator's intrinsic difficulty. Leave-one-generator-out validation, applied by excluding a generator's documents entirely rather than sampling around them, operationalizes structurally unseen generalization directly. A continuous paraphrase-strength sweep turns the unstable margin into a measured error curve rather than a binary judgment. Contrasting a quantification-corrected estimator against a naive one isolates how much of any observed gap is aggregation bias specifically. A capability-matched swap, which holds a fixed set of documents constant and varies only their calibration membership, is the only design that cleanly separates unseen-ness from intrinsic difficulty. A single joint statistical specification, rather than two independent tests, disentangles the generator-novelty and paraphrase-strength effects from one another. And a pre-committed criterion --- passing a cross-cell flatness test together with a per-cell floor, not merely a good average --- is what makes a generator-invariant-feature claim falsifiable rather than a post-hoc read of a table.

We evaluate this design on controlled human/synthetic mixtures built from a generator pool spanning multiple model families, using real human base text and established multi-generator machine-text collections, with paraphrased variants of the synthetic portion at several strength levels. In summary, we contribute:
\begin{itemize}[leftmargin=1.4em]
\item An identification and formalization of the asymmetry underlying corpus-level contamination estimation: the true mixture ratio is fixed by construction, while the observable statistic every estimator reads it off is free to vary with generator identity and paraphrase strength without bound --- an asymmetry that the deployed corpus-level prevalence estimators this paper stress-tests \citep{MonitoringAIContent2024} report a number from, but never bound.
\item A capability-matched seen/unseen swap that isolates genuine unseen-generator degradation from a held-out generator's intrinsic detectability --- a confound that the multi-generator robustness benchmarks nearest to this work \citep{RAID2024Benchmark,MAGE2024InTheWild} do not control for, because they compare \emph{across} generators rather than re-scoring identical documents under two calibration memberships.
\item A joint stress-disentanglement design and a quantification-correction contrast that separately attribute observed degradation to paraphrase strength, generator novelty, and aggregation bias rather than reporting one pooled error number.
\item A pre-committed, falsifiable cross-cell criterion for certifying (or ruling out) a generator-invariant statistical feature, together with an empirical account of which trajectory --- a fixable gap or a structurally obstructed one --- holds for current estimator designs.
\end{itemize}

\section{Problem Definition}

\noindent\textbf{Formal setup.} Consider a corpus $D$ built from a human base text set and a synthetic set produced by one or more generators, mixed at a total synthetic fraction $r \in [0,1]$ fixed by construction. An estimator $\hat{E}$ maps $D$ to an estimated fraction $\hat{r} = \hat{E}(D) \in [0,1]$. Let $G$ denote a pool of candidate text generators; for a held-out generator $g^* \in G$, the calibration set available to $\hat{E}$ is $G_{\text{cal}} = G \setminus \{g^*\}$, so that $D$'s synthetic documents produced by $g^*$ are structurally absent from calibration, not merely under-sampled within it. Paraphrase strength $k \in \{0,1,2,3\}$ is applied to the synthetic portion of $D$ before mixing, at increasing rewrite intensity. A candidate feature family $\phi$ has a per-condition human/synthetic separability $\mathrm{AUROC}_\phi(g,k)$. The object of study is the behavior of the absolute estimation error $\mathrm{AE} = |\hat{r} - r|$ as a joint function of $(g^*, k)$, and whether any $\phi$ retains stable separability across every tested $(g,k)$ pair.

\noindent\textbf{Why this is hard.} Because $G$ is open-ended and continually growing, no calibration set $G_{\text{cal}}$ can be exhaustive, so a meaningful robustness claim requires evaluating $\mathrm{AE}$ specifically when $g^*$ is structurally absent from $G_{\text{cal}}$, not merely when its documents are resampled. Because $k$ varies continuously, $\mathrm{AE}$ as a function of $k$ has no natural single-point summary --- the paraphrase margin is inherently a curve, not a threshold. Because $\hat{E}$ typically aggregates many per-document readings into one ratio, part of $\mathrm{AE}$ reflects a distribution-shift-induced aggregation bias rather than per-document classification error --- a bias documented experimentally in the quantification-learning literature for classify-and-count estimators under dataset shift \citep{QuantDatasetShift2023} --- and this bias is not removable by a fixed additive correction since it moves with $g^*$. Because $\mathrm{AE}(g^*, \cdot)$ alone conflates ``generator $g^*$ was unseen'' with ``generator $g^*$'s text is intrinsically hard to classify,'' attributing an observed gap to unseen-ness requires a design that holds $g^*$'s documents fixed and varies only calibration membership. Because $g^*$ and $k$ act on the same documents simultaneously, a univariate test of either factor risks misattributing part of one factor's effect to the other. And because a feature $\phi$ is only usefully ``invariant'' if $\mathrm{AUROC}_\phi(g,k)$ does not collapse in any tested condition, certifying invariance requires a worst-cell criterion evaluated over the full $(g,k)$ grid, not an average-case one.

\section{Method Overview}

At a conceptual level, this is a single measurement pipeline with six cooperating parts, organized around one design problem: the true synthetic fraction $r$ is fixed once a controlled corpus is built, while every candidate estimator's decision statistic is free to vary with generator identity and paraphrase strength, so the pipeline's job is to isolate exactly how much of that variation traces to each stress axis, to intrinsic per-generator difficulty, and to aggregation bias, separately.

Test corpora enter the pipeline built under a \textbf{leave-one-generator-out (LOGO) calibration} protocol, which excludes an entire generator's documents from calibration to operationalize structurally unseen generalization. The synthetic portion of each corpus passes through a \textbf{paraphrase-strength sweep}, rewriting it at increasing intensity so that estimator error can be traced as a continuous curve rather than a single before/after comparison. A \textbf{quantification-correction contrast} runs a bias-corrected estimator alongside a naive one to attribute part of any observed gap specifically to aggregation bias. A \textbf{capability-matched seen/unseen swap} holds a fixed set of test documents constant and varies only whether their generator was present at calibration, isolating genuine unseen-ness from a generator's intrinsic difficulty. A \textbf{joint stress-disentanglement design} fits generator novelty and paraphrase strength in one specification so that neither axis's estimated effect can silently absorb the other's. And a \textbf{generator-invariance probe} evaluates candidate statistical features across the full grid of held-out-generator and paraphrase-strength conditions to certify, or rule out, a feature that survives all of them at once.

\section{Methodology}

\subsection{Leave-One-Generator-Out Calibration}
\textbf{Challenge.} The space of possible text generators is open-ended and continually growing, so no calibration set can ever be exhaustive, and any robustness claim needs to distinguish generalization to a genuinely unseen generator from generalization to an unseen sample of an already-seen one.

\textbf{Insight.} Excluding an entire generator's documents from calibration --- rather than holding out a random sample that may still share the generator's stylistic fingerprint with training examples --- is what actually operationalizes ``structurally unseen'' generalization, at the cost of a smaller calibration pool on every fold.

\textbf{Method.} For each generator in the pool, all of its documents are removed from every estimator's calibration data, and that estimator is evaluated only on test corpora built from the excluded generator's held-out documents; this is repeated once per generator in the pool, alongside a matched in-distribution condition in which no generator is withheld, to give every leave-one-generator-out reading a same-harness baseline for comparison.

\subsection{Paraphrase-Strength Sweep}
\textbf{Challenge.} Paraphrasing continuously interpolates between clearly-synthetic and clearly-human text, so the separability margin --- and any single point-estimate of contamination --- is inherently unstable near that margin, independent of which estimator is used.

\textbf{Insight.} Representing paraphrase strength as an ordered, swept parameter rather than a binary toggle converts an unstable qualitative judgment into a quantitative error curve with an interpretable slope.

\textbf{Method.} The synthetic portion of each test corpus is paraphrased at strength levels $k \in \{0,1,2,3\}$, corresponding to increasing rewrite rounds of a controllable discourse-level paraphraser of the kind established for detector-evasion studies \citep{DipperParaphrase2023}, gated by a semantic-fidelity check (embedding-cosine similarity to the pre-paraphrase text) so that degenerate, meaning-destroying rewrites at high $k$ do not masquerade as a paraphrase effect; estimator error is measured at every level, always including the $k=0$ no-paraphrase control.

\textbf{Scope of this axis.} That paraphrasing defeats \emph{per-document} detectors is an established result \citep{DipperParaphrase2023,ReliablyDetectable2023}, not a hypothesis this paper tests; the open question is whether the effect survives aggregation into a proportion, where an averaging argument predicts partial cancellation. Accordingly the sweep here is deliberately \textbf{non-adversarial}: rewrite rounds are applied uniformly and independently to each synthetic document, with no coordination toward a target estimate and no budget allocated across the corpus. A cost-accounted, directionally coordinated adversary --- which could plausibly defeat the averaging argument that an independent one cannot --- is a distinct question about strategy-proofness and is deferred; the uncoordinated sweep measured here is the baseline any such adversarial result would have to be read against.

\subsection{Quantification-Correction Contrast}
\textbf{Challenge.} Aggregating many per-document readings into a single corpus-level ratio is a known source of systematic bias under distribution shift, and that bias moves with the held-out generator rather than washing out with a fixed calibration constant.

\textbf{Insight.} Because a quantification correction targets exactly the aggregate-proportion bias rather than the base classifier's per-instance error rate, comparing a corrected estimator's leave-one-generator-out gap against a naive estimator's gap under otherwise identical conditions attributes part of the total degradation specifically to aggregation bias.

\textbf{Method.} A naive classify-and-count estimator and a quantification-corrected estimator are run on identical test corpora across every leave-one-generator-out fold and paraphrase level, and their respective gaps between in-distribution and leave-one-generator-out error are compared directly. The correction itself is inherited, not invented: adjusted-count and prior-shift ratio estimators are standard instruments of the quantification-learning literature \citep{QuantDatasetShift2023,QuantPriorShift2018}, whose experimental treatment of classify-and-count bias under dataset shift is what licenses the challenge above. What that literature does \emph{not} cover is the shift studied here. Its shift model is a change in the label prior between training and test; the shift this paper induces is a change in the \emph{generative process behind one class} while the label prior is held fixed by construction. An off-the-shelf prior-shift correction is therefore a well-motivated candidate for closing the leave-one-generator-out gap, not a solution known in advance to apply --- which is precisely why the corrected-versus-naive difference is informative rather than decorative.

\subsection{Capability-Matched Seen/Unseen Swap}
\textbf{Challenge.} A leave-one-generator-out gap alone cannot distinguish ``this generator was unseen at calibration'' from ``this generator's text is intrinsically harder to detect regardless of calibration,'' since a comparison across different generators conflates the two explanations.

\textbf{Insight.} The only way to separate these two explanations is to hold a fixed set of test documents constant and vary only whether their generator was included in calibration --- any design that instead compares across different generators cannot make this separation cleanly.

\textbf{Method.} For a fixed set of documents produced by a given generator, estimator error is computed twice under otherwise identical conditions (pool size, compute budget, paraphrase strength): once with that generator included in calibration, and once with it excluded. The difference between the two is the unseen-ness effect, purged of any generator-specific difficulty confound.

\subsection{Joint Stress-Disentanglement Design}
\textbf{Challenge.} Generator novelty and paraphrase strength act on the same documents simultaneously, so testing either factor in isolation risks one axis's estimated effect silently absorbing part of the other's true effect.

\textbf{Insight.} A single joint statistical specification, rather than two independently run tests, is what lets both main effects be estimated net of one another, because both are fit from the same model rather than from two separately-fit models that could each misattribute variance to the wrong factor.

\textbf{Method.} Absolute estimation error is regressed jointly on calibration regime (in-distribution vs. leave-one-generator-out) and paraphrase strength $k$, with a random intercept per generator family to account for non-independence among generators that share architectural lineage, yielding two separately falsifiable estimands from one fit: the regime effect net of paraphrase strength, and the paraphrase-strength effect net of regime.

\subsection{Generator-Invariance Probe}
\textbf{Challenge.} A claim that a statistical feature is ``generator-invariant'' is only meaningful if it is checked in the worst tested condition, not on average --- a feature that separates well in most conditions and collapses in one has not been shown invariant in the sense that matters for practice.

\textbf{Insight.} Certifying invariance therefore requires a criterion that fails a feature the moment any single tested condition is bad, rather than a criterion that a strong average performance across conditions could satisfy despite a hidden failure cell. A second, easily-missed requirement follows: the criterion must demand \emph{evidence of} flatness rather than \emph{absence of evidence against} it, or a noisy, underpowered grid would certify invariance most readily --- exactly backwards.

\textbf{Method.} For each candidate feature family (reference-model log-likelihood/perplexity, stylometric n-gram and burstiness statistics, and sentence-embedding directions), human/synthetic separability (AUROC) is measured within every held-out-generator $\times$ paraphrase-strength cell. A feature is certified generator-invariant only if it simultaneously clears an absolute floor ($\mathrm{AUROC} \geq 0.70$) in every tested cell \textbf{and} passes a flatness criterion stated as an \textbf{equivalence test}: with a pre-committed equivalence margin $\delta_{\text{inv}}$ on the cross-cell AUROC spread, the feature is certified only when the upper confidence bound on that spread falls strictly \emph{inside} $\delta_{\text{inv}}$ (a two-one-sided-tests formulation), so that certification requires positive evidence of flatness. A conventional significance test of the null ``cross-cell variance $=0$'' is deliberately not used as the certifying leg: because such a test certifies on a failure to reject, both a thinner grid and a stricter multiplicity correction would make invariance \emph{easier} to declare, inverting the direction in which both should be conservative. Under the equivalence formulation the Holm correction applies to the equivalence rejections and therefore tightens certification, as intended. A feature whose cross-cell confidence bound is too wide to fall inside $\delta_{\text{inv}}$ but also does not demonstrably exceed it is scored \textbf{INCONCLUSIVE}, never certified and never reported as a refutation.

\section{Experiments}

\subsection{Experimental Setup}

\textbf{Datasets.} A human base corpus of real news and web prose provides the non-synthetic reference side, ensuring the pure-human self-consistency check ($\hat{r} \approx 0$ at $k=0$ with no generator withheld) is meaningful rather than trivially satisfied. Multi-generator machine-text collections spanning several model families \citep{RAID2024Benchmark,MAGE2024InTheWild} provide the synthetic side and cross-check corpus --- a pool of near-identical generators would make every leave-one-generator-out fold artificially easy, so family diversity in the pool is a design requirement, not a convenience. Paraphrased variants of the synthetic portion, produced by a controllable discourse-level rewriting pipeline \citep{DipperParaphrase2023} at strength levels $k \in \{0,1,2,3\}$, directly operationalize the paraphrase-strength axis. A human/model-comparison corpus with matched prompts on both the human and model side (HC3-style) supplies an additional, distributionally distinct human reference specifically for the pure-human self-consistency check, so that a passing check is not an artifact of one human source's domain-specific text statistics.

\textbf{Baselines} (sanity-check and confound-isolation comparators, not competitors to beat --- this is a measurement study, and no research question here is framed as ``our estimator beats baseline X''). A text-blind constant estimator, which always returns the calibration-set's prior mixture ratio while ignoring the corpus entirely, is the necessary ``why do this at all'' gate: if no real estimator beats this baseline in-distribution, there is no working estimator whose off-calibration degradation is worth measuring in the first place. The capability-matched seen/unseen swap (Methodology) is not a competing method but the confound-isolation arm for attributing degradation to unseen-ness specifically. The naive-versus-quantification-corrected contrast is a matched pair whose value lies in their difference, not in which one ``wins.''

\textbf{Metrics.} Absolute estimation error $\mathrm{AE} = |\hat{r} - r|$ and its condition-level mean; the leave-one-generator-out gap $\Delta_{\text{LOGO}}$ (mean LOGO error minus mean in-distribution error); the paraphrase-strength degradation slope; the seen-versus-unseen error difference on matched documents; per-feature AUROC and its cross-cell variance. Every significance test below is run per member of the tested estimator suite (M3(a)-(d)) or feature family and Holm-corrected within that suite before any RQ is scored, so a single favorable cell cannot be reported without accounting for the number of comparisons made.

\textbf{Power and inconclusive cells.} Every (generator, paraphrase-strength, mixture-ratio) cell is pre-declared to need $\geq 1000$ documents to detect the 0.03 absolute-error effect floor (A5); a cell realizing fewer documents, or whose bootstrap confidence interval on $\mathrm{AE}$ is wider than 0.03, is pre-declared \textbf{INCONCLUSIVE} before any hypothesis is scored on it --- never reported as a null result and never silently dropped from the denominator of cells tested. If a central test (RQ1's regime coefficient or RQ7's flatness-and-floor criterion) draws on an underpowered majority of its cells, that RQ is not scored PASS/FAIL at all; the claim narrows to the powered subset and the coverage gap is reported explicitly.

\textbf{Power at the generator level, which is the unit RQ1 actually tests.} The $\geq 1000$-document floor is a \emph{within-cell} requirement, and it is not the binding constraint on the headline test: because leave-one-generator-out confidence intervals are bootstrapped over generators rather than documents, RQ1's inferential unit is the generator, and the pool size caps that test's power independently of how many documents each cell holds. Stating the resulting bound explicitly, for a two-sided test at $\alpha = 0.05$ with 80\% power on the per-generator leave-one-generator-out gap: at the full pool of 12 generators the minimum detectable standardized effect is $d \approx 0.89$, so detecting the pre-declared 0.05 absolute-error gap requires the \textbf{between-generator standard deviation of that gap to be at most $\approx 0.056$}; at the walking skeleton's 3 generators the requirement tightens to $d \approx 3.26$ (SD $\lesssim 0.015$), which confirms the skeleton is a directional smoke test and cannot license any inference about the gap's magnitude. This between-generator dispersion is an assumption, not a measured quantity, and it is pre-declared as a gate in its own right: if the realized between-generator SD of the per-generator gap exceeds 0.056, RQ1 is reported as \textbf{INCONCLUSIVE at the tested pool size} --- with the realized SD and the pool size implied by it reported --- rather than as a null result, and the same rule applies to RQ4's paired swap and to RQ5's regime coefficient, both of which inherit the generator as their unit of replication.

\textbf{Research questions:}
\begin{enumerate}[leftmargin=1.6em]
\item RQ1 --- Does leave-one-generator-out validation produce higher error than in-distribution validation?
\item RQ2 --- Does estimation error increase monotonically with paraphrase strength?
\item RQ3 --- Does quantification correction reduce, without eliminating, the leave-one-generator-out gap?
\item RQ4 --- Is the leave-one-generator-out gap attributable to unseen-ness rather than a held-out generator's intrinsic difficulty?
\item RQ5 --- Does the calibration-regime effect remain significant once paraphrase strength is controlled for?
\item RQ6 --- Does the paraphrase-strength effect remain significant once calibration regime is controlled for?
\item RQ7 --- Does any candidate statistical feature certify as generator-invariant?
\end{enumerate}

\subsubsection{RQ1: Does leave-one-generator-out validation produce higher error than in-distribution validation?}
\textbf{Setup.} For every tested (estimator $\times$ generator family) cell, compute mean absolute error under in-distribution validation and under leave-one-generator-out validation, holding harness, scenario sampling, and paraphrase level ($k=0$) fixed between the two. Sample size: sized so that a 0.05 absolute-error difference is detectable with a generator-resampled bootstrap.

\textbf{Expected outcome.} Mean absolute error under leave-one-generator-out validation exceeds mean absolute error under in-distribution validation by at least 0.05, significantly, for the majority of tested estimators.

\textbf{Falsification condition.} Refuted if the leave-one-generator-out gap is not significantly greater than zero, or falls short of 0.05, for the majority of tested estimators --- that outcome would mean current estimator designs already generalize acceptably to structurally unseen generators without any further correction.

\textbf{Hypothesis (plain language).} Calibrating against a fixed generator pool does not, by itself, produce an estimator that stays reliable once the true generator is absent from that pool.

\textbf{Metrics.} Mean absolute error per condition; $\Delta_{\text{LOGO}}$ and its bootstrap confidence interval, resampled over generators rather than documents.

\textbf{Result.} \ph

\subsubsection{RQ2: Does estimation error increase monotonically with paraphrase strength?}
\textbf{Setup.} Within the in-distribution and leave-one-generator-out conditions alike, compare mean absolute error across paraphrase levels $k \in \{0,1,2,3\}$ via a monotonicity-oriented regression of error on $k$.

\textbf{Expected outcome.} Mean absolute error increases monotonically with paraphrase strength $k$, with a slope significantly greater than zero.

\textbf{Falsification condition.} Refuted if the estimated slope is not significantly positive, or if error is non-monotonic (e.g., peaks at an intermediate $k$ and declines) across the tested range --- either would mean paraphrase strength, as parameterized here, is not the driver of degradation it is hypothesized to be.

\textbf{Hypothesis (plain language).} Heavier paraphrasing of the synthetic portion makes it progressively harder for corpus-level estimators to recover the true mixture ratio.

\textbf{Why this is not a foregone conclusion.} That paraphrase degrades per-document detectors is established \citep{DipperParaphrase2023,ReliablyDetectable2023} and is not what this RQ tests. The open question is whether that degradation survives aggregation: if per-document errors induced by paraphrase were symmetric and independent across documents, they would partially cancel in a proportion, and the aggregate slope could be flat even while the underlying detector collapses. A significantly positive slope is therefore evidence that paraphrase-induced error is \emph{directional} at the document level, not merely large --- a distinct and non-obvious finding.

\textbf{Metrics.} Slope of mean absolute error on $k$, with confidence interval; per-level mean absolute error.

\textbf{Result.} \ph

\subsubsection{RQ3: Does quantification correction reduce, without eliminating, the leave-one-generator-out gap?}
\textbf{Setup.} Compute $\Delta_{\text{LOGO}}$ separately for the naive classify-and-count estimator and the quantification-corrected estimator on identical test corpora, and compare the two gaps via a paired bootstrap over generators.

\textbf{Expected outcome.} The corrected estimator's leave-one-generator-out gap is smaller than the naive estimator's gap by at least 30\%, but remains significantly greater than zero.

\textbf{Falsification condition.} Refuted toward ``correction is not load-bearing'' if the reduction in $\Delta_{\text{LOGO}}$ is below 30\% or not statistically distinguishable from zero; refuted toward ``correction fully resolves the gap'' if the corrected estimator's $\Delta_{\text{LOGO}}$ is not significantly greater than zero.

\textbf{Hypothesis (plain language).} Correcting for aggregation bias helps close part of the leave-one-generator-out gap, but the gap is not purely an aggregation-bias artifact.

\textbf{Metrics.} Paired difference in $\Delta_{\text{LOGO}}$ between corrected and naive estimators, with bootstrap confidence interval.

\textbf{Result.} \ph

\subsubsection{RQ4: Is the leave-one-generator-out gap attributable to unseen-ness rather than a held-out generator's intrinsic difficulty?}
\textbf{Setup.} For a fixed set of documents from a given generator, compute absolute error twice under matched pool size, compute budget, and paraphrase strength --- once with that generator included in calibration (seen), once excluded (unseen) --- and compare via a paired test over documents within generator.

\textbf{Expected outcome.} The seen-versus-unseen absolute-error difference on identical held-out documents is significantly greater than zero.

\textbf{Falsification condition.} Refuted if the seen-versus-unseen difference is not statistically distinguishable from zero --- that outcome would mean the leave-one-generator-out gap observed in RQ1 is driven by which generators happen to be difficult, not by calibration status, undermining any claim that ``unseen-ness'' specifically is the cause.

\textbf{Hypothesis (plain language).} Being unseen at calibration time degrades estimator performance on a generator's text beyond whatever intrinsic difficulty that generator's text already has.

\textbf{Metrics.} Paired seen-versus-unseen absolute-error difference, with confidence interval.

\textbf{Result.} \ph

\subsubsection{RQ5: Does the calibration-regime effect remain significant once paraphrase strength is controlled for?}
\textbf{Setup.} In the joint mixed-effects fit (absolute error regressed on calibration regime and paraphrase strength jointly, random intercept per generator family), read off the regime coefficient's significance after the paraphrase-strength term is included.

\textbf{Expected outcome.} The calibration-regime coefficient remains significant in the joint fit.

\textbf{Falsification condition.} Refuted if the regime coefficient is not significant once paraphrase strength is included jointly --- that would mean the apparent leave-one-generator-out effect observed univariately in RQ1 is actually attributable to a paraphrase-strength imbalance between conditions, not to calibration regime itself.

\textbf{Hypothesis (plain language).} Generator novelty degrades estimation independently of how heavily the text has been paraphrased, not merely as a byproduct of paraphrase-strength differences between conditions.

\textbf{Metrics.} Regime coefficient estimate, standard error, and significance from the joint fit.

\textbf{Result.} \ph

\subsubsection{RQ6: Does the paraphrase-strength effect remain significant once calibration regime is controlled for?}
\textbf{Setup.} In the same joint mixed-effects fit as RQ5, read off the paraphrase-strength coefficient's significance after the calibration-regime term is included.

\textbf{Expected outcome.} The paraphrase-strength coefficient remains significant in the joint fit.

\textbf{Falsification condition.} Refuted if the paraphrase-strength coefficient is not significant once calibration regime is included jointly --- that would mean the apparent paraphrase effect observed univariately in RQ2 is actually attributable to an imbalance in which generators appear at which paraphrase level, not to paraphrase strength itself.

\textbf{Hypothesis (plain language).} Paraphrase strength degrades estimation independently of whether the generator was seen at calibration, not merely as a byproduct of generator-novelty differences between conditions.

\textbf{Metrics.} Paraphrase-strength coefficient estimate, standard error, and significance from the joint fit.

\textbf{Result.} \ph

\subsubsection{RQ7: Does any candidate statistical feature certify as generator-invariant?}
\textbf{Setup.} For each candidate feature family, compute AUROC separability in every tested held-out-generator $\times$ paraphrase-strength cell, then apply the pre-committed flatness \textbf{equivalence} test --- a two-one-sided-tests procedure requiring the upper confidence bound on the cross-cell AUROC spread to fall inside the pre-committed margin $\delta_{\text{inv}}$ --- jointly with the 0.70 per-cell floor, Holm-corrected across the feature families tested. The equivalence formulation is load-bearing rather than stylistic: it makes certification require positive evidence of flatness, so that a thin grid or a strict multiplicity correction makes invariance \emph{harder} to declare rather than easier, which a conventional test of the null ``cross-cell variance $=0$'' would invert.

\textbf{Expected outcome.} No candidate feature family simultaneously passes the flatness equivalence test and the 0.70 floor across all tested cells --- the negative branch is the observed outcome, though a partial-invariance ranking across features remains reportable.

\textbf{Falsification condition.} Refuted (a generator-invariant feature exists) if some candidate feature's cross-cell spread is bounded inside $\delta_{\text{inv}}$ \emph{and} it clears the 0.70 floor in every tested cell --- that outcome would instead support the positive branch: identification of a candidate generator-invariant signature, reported together with its AUROC floor and its realized spread bound. A feature whose spread bound is too wide to fall inside $\delta_{\text{inv}}$ but that does not demonstrably exceed the margin either is scored \textbf{INCONCLUSIVE} and reported as such, since neither branch has been earned.

\textbf{Hypothesis (plain language).} No statistical feature tested here survives both an unseen generator and paraphrasing at once without its separability collapsing somewhere in the tested grid.

\textbf{Metrics.} Per-cell AUROC for each feature family; the cross-cell AUROC spread with its upper confidence bound, and the equivalence-test result against $\delta_{\text{inv}}$; minimum per-cell AUROC; the count of feature families scored INCONCLUSIVE, reported alongside those scored certified or refuted.

\textbf{Result.} \ph

\vspace{0.5em}
\noindent\textbf{Further evaluation.} With more time, a natural extension is validating whether the degradation law and feature-invariance verdict observed on controlled mixtures transfer to real, unlabeled corpora --- explicitly out of scope here because no independently verified ground-truth contamination fraction exists for real corpora, which is precisely why this study is confined to a controlled-mixture setting in the first place.

\section{Related Work}

Three prior works are closest, on three different dimensions, and each is conceded explicitly before the wider landscape is surveyed.

On the \emph{estimator} dimension --- what this paper stress-tests rather than competes with --- the nearest work is the distributional mixture framework for monitoring AI-modified content at scale \citep{MonitoringAIContent2024}, with its companion large-scale application to the scientific literature \citep{MappingLLMUse2024}. This is the canonical \emph{uncooperative} corpus-level proportion estimator: it fits a maximum-likelihood mixture against reference distributions to report the fraction of AI-modified content in a corpus, needing neither a watermark nor per-document labels, and it has been deployed on real review pools and paper collections whose curation consequences are non-hypothetical. We concede that it already does the thing --- estimate a corpus-wide proportion without generator cooperation --- and this paper proposes no replacement for it. We differentiate on what it never reports: its reference distributions are constructed from a fixed, known generator set and a fixed human anchor, and no result characterizes its error when the contaminating generator is absent from that set, when the text has been paraphrased, or how those two stresses interact. This paper's object of study is precisely that unreported error surface, for this class of estimator.

On the \emph{stress-axis} dimension, the nearest works are the multi-generator robustness benchmarks for machine-generated-text detection: a shared benchmark evaluating detector robustness across many generator families crossed with adversarial attacks including paraphrase \citep{RAID2024Benchmark}, and an out-of-distribution testbed spanning unseen generators and unseen domains \citep{MAGE2024InTheWild}. We concede this dimension outright and without hedging: these benchmarks already occupy both of this paper's stress axes, in text, at a scale of generator coverage this study does not attempt to match, and they are the reason the per-document half of the story requires no re-litigation here. The delta is the estimand and the identification design, not the axes. First, both score a per-document decision; neither reports the error of a corpus-wide \emph{proportion}, and the two levels are not interchangeable, since aggregation can either cancel per-document error (the averaging argument) or systematically bias it (the quantification argument). Second, neither runs a quantification correction, so neither can attribute any part of an observed degradation to aggregation bias specifically. Third, and most consequentially for the causal claim, both compare performance \emph{across} different generators, which conflates ``this generator was unseen at calibration'' with ``this generator's text is intrinsically hard'' --- the confound the capability-matched swap in this paper's Methodology is built to remove by re-scoring identical documents under two calibration memberships.

On the \emph{estimand-theory} dimension, the nearest work is a sample-complexity analysis for watermark-proportion estimation \citep{SampleComplexitiesofEsti700224}, which targets the identical quantity this paper does --- a mixing proportion of machine-generated text in an observed sample --- with provably near-optimal sample complexity, more rigorous on the estimator dimension than anything attempted here. We concede that dimension outright: its estimators are more sample-efficient than any of ours \emph{when their assumption holds}. We differentiate on the assumption itself --- its efficiency guarantees hold only under a cooperative watermarking scheme (a known pseudorandom green-list rule), and it reports no result for a generator absent from that scheme and no paraphrase-attack condition, so it does not touch this paper's actual question: whether \emph{any} distributional estimator remains reliable when the generator is uncooperative and structurally unseen, and under paraphrase. This paper is positioned as the uncooperative/unseen-generator complement to that estimator theory, not a replacement for it.

Taken together: measuring a corpus-wide proportion's degradation rather than a per-document decision's, under a structurally withheld generator crossed with paraphrase strength, with unseen-ness separated from intrinsic difficulty and aggregation bias separated from classification error, is what none of the three nearest works does and what this paper is built to do.

\textbf{Corpus-level prevalence estimation without generator cooperation.} Beyond the framework conceded above, a growing line of work reports population-level AI-text prevalence rather than per-document verdicts, and treats the resulting number as an input to a monitoring or curation decision \citep{MonitoringAIContent2024,MappingLLMUse2024}. What unites these studies methodologically is also what makes them this paper's target: each validates its estimator against generators it selected, on a human reference it fixed in advance, and reports a point estimate with sampling uncertainty but no term for calibration-set misspecification. The failure this paper measures is not that these estimates are wrong on their reported benchmarks; it is that the reported uncertainty does not include the dominant term when the true generator is one the auditor never anticipated.

\textbf{Multi-generator robustness benchmarks for text detectors.} RAID \citep{RAID2024Benchmark} evaluates detectors across a large cross of generator families, decoding strategies, and adversarial attacks --- paraphrase among them --- and MAGE \citep{MAGE2024InTheWild} constructs explicit out-of-distribution splits over unseen generators and unseen domains. Both establish, for text and not by analogy, that detector performance degrades substantially off the generators it was built against; this study takes that as settled premise rather than as a finding to reproduce. Their reported quantity throughout is per-document accuracy or AUROC, aggregated as an average over conditions rather than certified at the worst condition, which leaves two things open that this paper closes: the behavior of the aggregate proportion itself, and a worst-cell criterion for whether any feature survives every condition at once. Both also serve as this study's instrument base --- their multi-family generator pools are what make a leave-one-generator-out design non-trivial, since a pool of near-identical generators would render every fold artificially easy.

\textbf{Paraphrase as a detector stressor.} That a bounded, quality-preserving paraphrase defeats state-of-the-art per-document detectors is established: a discourse-level paraphraser trained for this purpose evades zero-shot, watermark, and trained detectors alike, with retrieval over previously-generated text offered as the defense \citep{DipperParaphrase2023}, and a companion analysis argues the vulnerability is general to the detection setting rather than specific to one detector \citep{ReliablyDetectable2023}. This paper adopts that paraphraser class as its rewrite instrument and adopts the result as premise. The unresolved question it inherits is one level up and is not answered by either work: both report per-document evasion at a single attack setting, whereas the quantity a curation decision consumes is an aggregate proportion, over which independent per-document errors might partially cancel, and the relevant object is an ordered strength curve rather than a before/after pair.

\textbf{Quantification under distribution shift.} The correction contrasted against naive classify-and-count in this study is inherited from the quantification-learning literature, which studies exactly the problem of estimating class prevalence rather than labeling instances. An experimental investigation of binary quantification under dataset shift documents that classify-and-count acquires systematic aggregate bias when the test distribution departs from the training one, and characterizes when adjusted-count corrections recover \citep{QuantDatasetShift2023}; a formal treatment derives the prior-shift ratio estimator and its extensions \citep{QuantPriorShift2018}. This literature is the source of both the instrument and the challenge motivating it, and is conceded as such. Its shift model, however, is a change in the \emph{label prior} between training and test. The shift this paper induces is different in kind: the label prior $r$ is fixed by construction while the generative process behind one class is replaced by an unseen one, so whether a prior-shift correction transfers is an open empirical question here rather than a settled guarantee --- which is what makes the corrected-versus-naive gap difference a measurement rather than a formality.

\textbf{Per-document machine-generated-text detection.} Zero-shot detectors score a candidate passage against the local curvature of a source model's log-probability \citep{DetectGPT2023}, requiring likelihood access (or a close surrogate) to the generating model. Rewrite-distance detectors instead learn a distance between a text and its own paraphrase \citep{LearnToDistance2023}, motivated explicitly by the observation that curvature-based signals ``depend strongly on assumptions'' that break under distribution shift. Recovery-based detectors avoid likelihood access altogether by disrupting and then recovering a passage with a single black-box call \citep{DR2023RecoveryDetect}, scoring recovery similarity as a proxy for posterior concentration. A calibration layer wraps such detectors in a pairwise Markov-random-field model over token-level scores before aggregating to a passage-level decision \citep{MarkovCalibration2023}. All four output a per-document decision or score, and none reports a corpus-level proportion; robustness across generators and paraphrase for detectors of exactly this kind is what the benchmarks conceded above \citep{RAID2024Benchmark,MAGE2024InTheWild} already measure. The aggregation step the calibration layer treats as a noisy preprocessing detail is exactly the step this paper's quantification-correction contrast and leave-one-generator-out design make into the object of study.

\textbf{Provenance via generation-time cooperation.} A watermarking scheme embeds a detectable statistical signal at generation time by pseudorandomly biasing the sampled vocabulary \citep{KirchenbauerWatermark2023}, reducing detection to a one-proportion test with no model access required at detection time --- but only for generators that implement the scheme. Its own motivation concedes that non-watermark, post-hoc detectors degrade rapidly and fail to generalize across model generations and sampling strategies, which is the failure this paper measures directly, but a watermarking-based approach cannot help estimate contamination from a generator the auditor never controlled or instrumented in the first place --- the complement, not a competitor, to the corpus-auditing setting studied here. The sample-complexity analysis conceded above \citep{SampleComplexitiesofEsti700224} formalizes detection under exactly this kind of cooperative scheme, sharpening the same complement relationship at the level of a proven estimator rather than an empirical one.

\textbf{Detector generalization to unseen generators in other modalities.} A contrastive-training approach improves a diffusion-image detector's accuracy on generator variants unseen during training by augmenting with contrastively reconstructed hard samples \citep{DRCT2024DiffusionRecon}, and its own motivation states that detectors trained on a narrow set of models exhibit severe degradation against unseen diffusion variants. This is a cross-modal precedent for the failure mode rather than this paper's nearest neighbor --- that role belongs to the text benchmarks conceded above --- and it is retained here for a specific reason: it is the clearest existing demonstration that the unseen-generator problem is a property of the detection setting rather than an artifact of any one modality's statistics, which is part of why an estimator-level treatment is warranted rather than a detector-level patch. Because it fixes the problem via training-time augmentation of a binary detector rather than characterizing a corpus-level estimator's degradation rate, and because it has no paraphrase-strength axis to cross with generator novelty, it leaves the measurement this paper performs untouched.

\textbf{Synthetic-data mixing theory and adjacent contamination measurement.} A multi-entity recursive-training study models several systems mixing reused real data with generation-indexed synthetic data from all entities \citep{RecursiveMultiEntity2024}, quantifying convergence toward homogeneity under a controlled synthetic-mixing fraction; a companion theoretical analysis derives the exact conditions on that mixing parameter under which iterative retraining on the resulting corpus is a stable contraction rather than a collapsing process \citep{IterativeRetrainStability2024}. Both treat the mixing fraction as a controlled input to study its downstream training-dynamics consequences, establishing that the fraction is the causally load-bearing quantity --- but neither provides an inverse estimator that recovers an unknown fraction from an observed corpus, let alone one whose robustness to unseen generators is characterized. A related contamination-measurement study documents preference leakage in LLM-as-judge pipelines \citep{PreferenceLeakage2024}, measuring the effect for specific, related generator and judge families by design --- an adjacent contamination-measurement result bound to the model families it studies, rather than a leave-one-generator-out or paraphrase-strength stress test of a corpus-level estimator.

\section{Conclusion}

This work reframes an assumption implicit in every corpus-level contamination estimator --- that reliability demonstrated against a fixed calibration pool of generators carries over to generators never seen at calibration time --- as a measurable degradation law with a precise identifiability structure, rather than an unexamined default. The central methodological wager is that a leave-one-generator-out gap is only interpretable once it has been separated from a held-out generator's intrinsic difficulty, disentangled from paraphrase strength, and decomposed into an aggregation-bias component and a residual component; a plain before/after comparison collapses all of these into one number that could mean several different things. Whichever regime dominates for current estimator designs, the finding is decision-relevant: if quantification correction closes most of the gap, off-calibration robustness is an engineering problem worth investing in; if a generator-invariant feature is certified, a concrete fix exists; if neither holds, teams relying on a single point estimate of contamination to make filtering decisions are relying on a number whose error is structurally unbounded exactly when the generator responsible for contamination is the one they didn't anticipate. Six limitations bound this claim from the outset and are natural next steps rather than afterthoughts, not generic caveats. \textbf{External validity to real corpora} is bounded, not demonstrated: validation is confined to controlled mixtures with a researcher-set ground truth, since no independently verified contamination fraction exists for real, unlabeled corpora. \textbf{Construct validity of ``unseen'' itself} depends on held-out generators being structurally distinct from the calibration pool rather than a same-family resample that leaks a shared architectural fingerprint; this is checked via a per-generator-family random intercept and generator-resampled confidence intervals, but that check is only adequately powered once the pool is scaled beyond the walking-skeleton pilot's 3 generators, so any pilot-stage attribution of a gap specifically to unseen-ness (rather than to family-shared artifacts) is provisional until the scaled check runs. \textbf{Pool representativeness} is a design choice, not a sampled population: a pool that turns out too homogeneous across the tested generators would trivially shrink the leave-one-generator-out gap for reasons unrelated to the phenomenon under study, so the degradation-rate estimate is scoped to the tested generator families and does not claim to generalize to families outside the pool without further evidence. \textbf{The invariance verdict is three-valued, not two-valued}: because the flatness criterion is stated as an equivalence test against a pre-committed margin rather than as a failure to reject a variance null, a feature can end the study certified, refuted, or INCONCLUSIVE, and the third outcome is a statement about the tested grid's resolution rather than about the feature --- reported as such rather than folded into the negative branch. \textbf{The paraphrase axis is non-adversarial by construction}: rewrite strength is applied uniformly and independently, with no coordination toward a target estimate and no budget allocated across the corpus, so the degradation law reported here bounds what an uncoordinated rewriting process does and says nothing about what a strategic adversary optimizing against the estimator could achieve --- a separate question about strategy-proofness that this protocol is designed to be the baseline for, not to answer. Finally, the study is scoped to one language and one text domain, leaving code, multilingual text, and multimodal contamination for future work built on the same protocol.

\vspace{1em}
\hrule
\vspace{1em}

\section*{Pre-handoff checklist}
\begin{itemize}[leftmargin=1.4em]
\item \textbf{Organization/navigation} --- narrative flows problem $\to$ method $\to$ experiments $\to$ related work without forward references to undefined terms.
\item \textbf{Terminology discipline} --- ``contamination fraction,'' ``leave-one-generator-out,'' ``paraphrase strength,'' ``generator-invariant feature'' each defined at first use in Problem Definition, used consistently thereafter; no internal stage IDs (\texttt{sp\_*}, \texttt{m\_*}, \texttt{rq\_*}) leaked into prose --- checked directly against the discussion pipeline log.
\item \textbf{Notation with worked example} --- the formal setup section defines $(D, r, \hat{E}, \hat{r}, G, g^*, G_{\text{cal}}, k, \phi)$ and the Problem Definition/Methodology sections consistently reuse this notation; no separate worked numeric example is included since no numbers exist yet (would require fabricating results).
\item \textbf{RQ framing matches outcomes} --- all 7 RQs are graded/quantitative (error differences, slopes, coefficients, AUROC floors), each answered in its own subsection with an explicit falsifiable expected outcome and a distinct falsification condition, not a binary yes/no.
\item \textbf{No typesetting artifacts} --- only explicit placeholder markers in Results, per the drafting convention this is a pre-registered spec, not a report.
\item \textbf{Artifact plan} --- the controlled-mixture test suite, paraphrase pipeline configuration, and estimator implementations are the natural release artifacts; not yet built, so not claimed as available.
\end{itemize}

\section*{Hand-off}
This is a draft, not a finished paper --- experiments have not run. Natural next steps: prose polish and venue-specific LaTeX formatting once a target venue is chosen, or revising an earlier stage if piloting surfaces a problem (e.g., if the leave-one-generator-out gap in RQ1 is much smaller than expected).

\bibliographystyle{plain}
\bibliography{refs}

\end{document}
